% titel: Blatt 5, A5: Vergleich des Atomradius
% skript: Mortimer Kap. 7 (chemie: Periodizität, Atomradius)

Der Atomradius eines Elements lässt sich relativ zu Nachbarelementen abschätzen:
in einer Periode nimmt er nach rechts ab (grössere $Z_{\mathrm{eff}}$), in einer
Gruppe nach unten zu (mehr Schalen).

\begin{werkzeugkasten}[Radius-Vergleich zweier Atome]
\begin{itemize}
  \item Gleiche Periode $\to$ das weiter links stehende Atom hat den grösseren
        Radius (geringere Kernladung).
  \item Gleiche Gruppe $\to$ das weiter unten stehende Atom hat den grösseren
        Radius (zusätzliche Schale).
  \item Bei verschiedenen Gruppen und Perioden: Schalenzahl hat meist den
        grösseren Einfluss $\to$ das Element mit der höheren Periode ist grösser.
\end{itemize}
\textit{Skript:} Mortimer Kap. 7 (Periodizität).
\end{werkzeugkasten}

\begin{achtung}[Darauf achten]
\begin{itemize}
  \item Radiusvergleiche nur bei \textit{Atomen} (nicht Ionen) vornehmen.
  \item Reihenfolge im PSE auf den Punkt genau vergleichen --- bei Paaren wie
        P/Cl stehen beide in Periode 3, links hat P den grösseren Radius.
  \item P und As: gleiche Gruppe, As liegt eine Periode tiefer und ist grösser.
\end{itemize}
\end{achtung}
